% GENERATED FILE. Edit canonical lessons, then run npm run generate:books.
% canonical-source-hash: fnv1a64:57807a5dc9562dc9
% canonical-lessons: LA-C289-si, LA-C289-si-pluit, LA-C289-pendet-ex, LA-R289-if

\chapter{If}
\label{ch:la-if}
\hlchaptermodality{289}

\begin{chapteropening}
\textbf{By the end of this chapter:} \emph{I can make a real condition with sī and nisi, and say that something depends on something else.}

Complete LA-R289-if: I can make a real condition with sī and nisi, and say that something depends on something else.
\end{chapteropening}

\section[sī, nisi]{sī — if}

\label{lesson:LA-C289-si}



\begin{quote}
Your very first politeness gave you \textbf{sīs}, a squeezed \emph{sī vīs}, \textquotedblleft{}if you wish\textquotedblright{}. And you own \textbf{nisi}, \textquotedblleft{}unless\textquotedblright{}.

Both have the word for \textquotedblleft{}if\textquotedblright{} inside them. Now it stands on its own.
\end{quote}

\subsection*{You'll want to know: sī}
\textbf{sī} — \textquotedblleft{}if\textquotedblright{}. It opens the condition, and the condition usually comes first:

\begin{quote}
\textbf{Sī vīs, ambulāmus.} — If you like, we'll go for a walk. \textbf{Sī pluit, nōn ambulō.} — If it rains, I don't go walking.
\end{quote}

\textbf{nisi}, \textquotedblleft{}unless, if not\textquotedblright{}, opens the same slot the other way round:

\begin{quote}
\textbf{Nisi pluit, ambulō.} — Unless it rains, I go walking.
\end{quote}

\emph{Sī nōn} also exists, and it puts the weight on the \emph{nōn}. \emph{Nisi} is the everyday way to say \textquotedblleft{}unless\textquotedblright{}.

\begin{cousinweb}
\textbf{sī} is inherited, usually traced to a small Proto-Indo-European word \emph{*swe} (EtymDB 2.1); Lewis \& Short derive it from a pronoun stem and record an older spelling, \emph{seī}. The details are not settled. \textbf{nisi}, as its own lesson said, is \emph{ne}, \textquotedblleft{}not\textquotedblright{}, in front of \emph{sī}.
\end{cousinweb}

\subsection*{Your turn}
Say these aloud:

\begin{itemize}
\raggedright
  \item \emph{sī}, then \emph{nisi}
  \item \emph{Sī vīs, ambulāmus.}
  \item \emph{Sī pluit, nōn ambulō.}
  \item \emph{Nisi pluit, ambulō.}
  \item \emph{Sī vīs, iter Rōmam facimus.}
\end{itemize}

\begin{tcolorbox}[breakable,colback=teal!4,colframe=teal!35!black,title={Before you move on}]
Say \textquotedblleft{}if\textquotedblright{}. (\emph{\textbf{Sī.}}) Say \textquotedblleft{}unless\textquotedblright{}. (\emph{\textbf{Nisi.}}) What is \emph{sīs} squeezed from? (\emph{\textbf{Sī vīs}}, if you wish.) And where does the \emph{sī} clause usually stand? (\textbf{First.})
\end{tcolorbox}
\section[sī pluit, in vīllā manēmus]{sī pluit — a condition that may well be true}

\label{lesson:LA-C289-si-pluit}



\begin{quote}
Say \textquotedblleft{}if it rains, I don't go walking\textquotedblright{}. (\emph{Sī pluit, nōn ambulō.})

That is a \textbf{real} condition: it may well rain, and if it does, the second half follows. Latin keeps both halves plain.
\end{quote}

\begin{grammarlens}[title={Grammar Lens: ordinary verbs on both sides}]
\begin{quote}
\textbf{Sī pluit, in vīllā manēmus.} — If it rains, we stay in the farmhouse.
\end{quote}

Each half has the tense the facts need. A real condition about the past takes the perfect, the way any past event does:

\begin{quote}
\textbf{Sī herī in forō fuistī, amīcum vīdistī.} — If you were in the forum yesterday, you saw my friend.
\end{quote}

The speaker does not know whether you were there. If you were, the seeing follows. Nothing in the verbs is pretending.
\end{grammarlens}

\begin{culture}
The second half can be a command, or the \emph{dēbēs} of advice:

\begin{quote}
\textbf{Sī fessus es, quiēsce.} — If you are tired, rest. \textbf{Sī aeger es, quiēscere dēbēs.} — If you are ill, you should rest.
\end{quote}

\emph{Quiēsce} is the command of \emph{quiēscō}: the stem with \emph{-e}, like \emph{ignōsce} in \emph{ignōsce mihi}.
\end{culture}

\subsection*{Your turn}
Say these aloud:

\begin{itemize}
\raggedright
  \item \emph{Sī pluit, in vīllā manēmus.}
  \item \emph{Sī herī in forō fuistī, amīcum vīdistī.}
  \item \emph{Sī fessus es, quiēsce.}
  \item \emph{Sī aeger es, quiēscere dēbēs.}
  \item a real condition about tomorrow's weather, and what you do if it comes true
\end{itemize}

\begin{tcolorbox}[breakable,colback=teal!4,colframe=teal!35!black,title={Before you move on}]
Say \textquotedblleft{}if it rains, we stay in the farmhouse\textquotedblright{}. (\emph{\textbf{Sī pluit, in vīllā manēmus.}}) Which tense goes after \emph{sī} for a real past? (The \textbf{perfect}, like any event.) Say \textquotedblleft{}if you are tired, rest\textquotedblright{}. (\emph{\textbf{Sī fessus es, quiēsce.}})
\end{tcolorbox}
\section[pendet ex]{pendet ex — it depends}

\label{lesson:LA-C289-pendet-ex}



\begin{quote}
You own \textbf{pendeō}, \textquotedblleft{}to hang\textquotedblright{}. Say \textquotedblleft{}if it rains, we stay in the farmhouse\textquotedblright{}. (\emph{Sī pluit, in vīllā manēmus.})

Asked whether you will walk tomorrow, the honest answer is often neither yes nor no.
\end{quote}

\subsection*{You'll want to know: pendet ex}
\textbf{pendet ex} — \textquotedblleft{}it depends on\textquotedblright{}, word for word \textquotedblleft{}it hangs from\textquotedblright{}. \textbf{ex} is \textquotedblleft{}out of, from\textquotedblright{}, and the thing it hangs from goes in the ablative:

\begin{quote}
\textbf{Tōtum ex tempestāte pendet.} — It all depends on the weather.
\end{quote}

Cicero says it of hope: \textbf{spēs pendet ex fortūnā}, hope hangs on fortune (Lewis \& Short quote it under \emph{pendeo}). English \textquotedblleft{}depend\textquotedblright{} is the same picture, \emph{dē-pendēre}, to hang down from.

\subsection*{You'll want to know: in tuā potestāte est}
When it depends on the person you are talking to:

\begin{quote}
\textbf{In tuā potestāte est.} — It is up to you; it is in your power.
\end{quote}

You own \textbf{potestās}, \textquotedblleft{}power, the ability to act\textquotedblright{}; \emph{in tuā potestāte} is \textquotedblleft{}in your power\textquotedblright{}. Cicero, writing about fate, asks what is \emph{in nostrā potestāte}, in our power.

\begin{cousinweb}
\textbf{pendeō}, as its own lesson said, is the intransitive partner of \emph{pendō}, \textquotedblleft{}to weigh, to hang up\textquotedblright{}, usually traced to Proto-Indo-European \emph{*(s)pend-}, \textquotedblleft{}to stretch, to pull\textquotedblright{} (EtymDB 2.1). \textbf{ex} is inherited from Proto-Indo-European \emph{*h1eghs}, \textquotedblleft{}out\textquotedblright{} (EtymDB 2.1), and English \emph{exit} begins with it. A thing that \emph{pendet ex} is held up by something else, and falls if that gives way.
\end{cousinweb}

\subsection*{Your turn}
Say these aloud:

\begin{itemize}
\raggedright
  \item \emph{pendet ex}
  \item \emph{Tōtum ex tempestāte pendet.}
  \item \emph{Spēs pendet ex fortūnā.}
  \item \emph{In tuā potestāte est.}
  \item \emph{Sī pluit, in vīllā manēmus. Tōtum ex tempestāte pendet.}
\end{itemize}

\begin{tcolorbox}[breakable,colback=teal!4,colframe=teal!35!black,title={Before you move on}]
Say \textquotedblleft{}it all depends on the weather\textquotedblright{}. (\emph{\textbf{Tōtum ex tempestāte pendet.}}) Say \textquotedblleft{}it is up to you\textquotedblright{}. (\emph{\textbf{In tuā potestāte est.}}) What does \emph{pendet ex} say, word for word? (\textbf{It hangs from.})
\end{tcolorbox}
\section[(review)]{Review — if, and what it depends on}

\label{lesson:LA-R289-if}



\begin{quote}
Say \textquotedblleft{}if you like, we'll go for a walk\textquotedblright{}. (\emph{Sī vīs, ambulāmus.}) Then \textquotedblleft{}unless it rains\textquotedblright{}. (\emph{Nisi pluit.})
\end{quote}

\begin{grammarlens}[title={What you've built}]
Three lessons, and the real condition is yours. \emph{Sī} or \emph{nisi} first, then the result, and ordinary verbs on both sides: present, perfect, or a command. And when the answer is neither yes nor no: \emph{tōtum ex tempestāte pendet}, or, handing it back, \emph{in tuā potestāte est}.

\begin{quote}
\emph{— Crās iter Rōmam facimus?} \emph{— Ex tempestāte pendet. Sī pluit, in vīllā manēmus; nisi pluit, iter facimus.}
\end{quote}

\emph{Crās}, \textquotedblleft{}tomorrow\textquotedblright{}, with a present verb: Latin speech often lets the present stand for a plan.
\end{grammarlens}

\subsection*{Your turn}
Say these aloud:

\begin{itemize}
\raggedright
  \item \emph{Ex tempestāte pendet.}
  \item \emph{Sī pluit, in vīllā manēmus; nisi pluit, iter facimus.}
  \item \emph{Sī abhinc trēs annōs Rōmae fuistī, forum vīdistī.}
  \item \emph{In tuā potestāte est.}
  \item a plan of your own for tomorrow, with \emph{sī}, \emph{nisi}, and what it depends on
\end{itemize}

\begin{tcolorbox}[breakable,colback=teal!4,colframe=teal!35!black,title={Before you move on}]
Say \textquotedblleft{}it depends on the weather\textquotedblright{}. (\emph{\textbf{Ex tempestāte pendet.}}) Say \textquotedblleft{}it is up to you\textquotedblright{}. (\emph{\textbf{In tuā potestāte est.}}) Say \textquotedblleft{}unless\textquotedblright{}. (\emph{\textbf{Nisi.}}) And from the last chapter, say \textquotedblleft{}better\textquotedblright{}. (\emph{\textbf{Melior.}})
\end{tcolorbox}
